% !TEX root = ../../Extensions of IQC Duality to Partial Integral Equations.tex
Theorem~\ref{thm:dual-KYP} provides a robust-performance certificate for PIE
systems in terms of the dual IQC. For analysis, this merely reformulates the
primal certificate, since primal--dual equivalence preserves feasibility. The
benefit of the dual formulation therefore appears in controller
synthesis. Both controller constructions below use the same synthesis LPI,
obtained by applying \eqref{eq:dual-kyp} to the corresponding closed-loop
PIE. The choice of dual uncertainty filter changes only the realization of the
controller recovered from this LPI. Related applications of IQC duality
include the robust-estimation result of Scherer and
K\"ose~\cite{scherer2008robustness}, and the disturbance-feedforward synthesis
results of K\"ose and Scherer~\cite{kose2009robust} and Venkataraman and
Seiler~\cite{venkataraman}.

\dualcontrollerfigure

\subsection{Primal and Dual Systems}
Let $G$, with zero initial condition $\mbf{x}(0)=0$, have input $w$,
control input $u$, and output $z$. Its primal and dual PIE state-space
realizations are summarized in the corresponding system-component boxes.
If $\Theta\neq I$, assume that the signals entering its dynamic part are
available.

\subsection{State Feedback LPI}

We now state the main synthesis condition. On the dual realization, the change
of variables represented by $\mcl Z$ makes the controller-dependent terms
affine in the decision operators. Feasibility therefore returns both a
robust-performance certificate and the gains of the primal controller.




\begin{corollary}\label{thm:state-feedback}
    Let $G$ be the controlled PIE plant defined in
    Eq.~\eqref{eq:primal-system-components}.
    Let $\Delta\in\mbf{\Delta}$ and exist a primal (hard) IQC $(\Psi,\mcl{V})$ that has the hard
    $(\Theta,J_{\mbf n_{\mrm{z}},\mbf n_{\mrm{w}}})$ factorization. Assume that
    \begin{itemize}
        \item The interconnection $K\star P\star\Delta$ is well-posed for all $\Delta\in\mbf{\Delta}$.
        \item The IQC in Definition~\ref{def:uncertainty-iqc} is satisfied for all $\Delta\in\mbf{\Delta}$.
    \end{itemize}
     Moreover there exist, $\mcl{Z},\mcl{P}\in\Pi_4$ with $\mcl{P} = \mcl{P}^* \succ 0$, such that the following statement holds
    \begin{equation}\label{eq:controller-kyp}
        \begin{gathered}
            \bmat{
                \underline{\mcl T}^*\mcl P\underline{\mcl A}_0
                +
                \underline{\mcl T}^*\mcl Z\bmat{\mcl B_{\mrm{u}}^*&0}
                +(*)
                &
                \underline{\mcl T}^*\mcl P\underline{\mcl B}_0
                +
                \underline{\mcl T}^*\mcl Z\mcl D_{\mrm{zu}}^*
                \\
                (*)&
                \underline\epsilon I
            }\\
                +\bmat{\underline{\mcl C}^*\\\underline{\mcl D}^*}
                J_{\mbf n_{\mrm{w}},\mbf n_{\mrm{z}}}
                \bmat{\underline{\mcl C}&\underline{\mcl D}}
                \preceq0.
        \end{gathered}
\end{equation}
where
\begin{equation}\label{eq:feedback-system}
\mbf{D}(\Theta)\bmat{0\star G^\top\\I}\underline{w}
:=
\left\{
\bmat{
\underline{\mcl T}\dot{\underline{\mbf{\xi}}}(t)\\
{\underline{v}}(t)
}
=
\bmat{
\underline{\mcl A}_0 & \underline{\mcl B}_0\\
\underline{\mcl C} & \underline{\mcl D}
}
\bmat{
\underline{\mbf{\xi}}(t)\\
\underline{w}(t)
}
\right.
\end{equation}
and recover the controller gains from
\[
\bmat{\mcl K_{\mrm{G}}&\mcl K_\Theta}=\mcl Z^*\mcl P^{-1}.
\]
Assume that $\mcl D_{21}\mcl D_{\mrm{G}}+\mcl D_{22}$ is invertible, with
\[
(\mcl D_{21}\mcl D_{\mrm{G}}+\mcl D_{22})^{-1}\in\Pi_4
\]
Then the interconnection $K\star P\star\Delta$ with $\mbf{x}(0)=0$ is robustly stable and robust performance is achieved.
\end{corollary}
\begin{proof}
    See \ref{proof:state-feedback}.
\end{proof}

The operator $\mcl P$ certifies the dual closed loop, while $\mcl Z$ contains
the transformed controller gains. The recovery formula maps these decision
operators back to the controller acting on the primal plant and filter states.
The invertibility assumption ensures that the resulting algebraic
interconnection is well-posed.


\subsection{\texorpdfstring{$\underline{\Psi}_\Delta = I$}
{Dual-filter identity} controller realization}

When $\Psi_\Delta=I$ and $\underline{\Psi}_\Delta=I$, no multiplier state is
required in the controller. For a state-feedback operator $\mcl{K}_{\mrm{G}}$, set
$u(t)=\mcl{K}_{\mrm{G}}\mbf{x}_{\mrm{G}}(t)$. The resulting closed-loop PIE
$K\star G$ is
\begin{equation*}
\bmat{\partial_t(\mcl{T}_{\mrm{G}}\mbf{x}_{\mrm{G}})(t)\\z(t)}
=
\bmat{
\mcl{A}_{\mrm{G}}+\mcl{B}_{\mrm{u}}\mcl{K}_{\mrm{G}}&\mcl{B}_{\mrm{G}}\\
\mcl{C}_{\mrm{G}}+\mcl D_{\mrm{zu}}\mcl{K}_{\mrm{G}}&\mcl{D}_{\mrm{G}}
}
\bmat{\mbf{x}_{\mrm{G}}(t)\\w(t)}.
\end{equation*}
Thus the identity-filter case requires no additional controller state. The
controller is memoryless in time and acts directly on the measured PIE state.

\subsection{Dynamic-filter controller realization}

Dynamic-IQC feedback constructions with additional controller memory have
also been studied for uncertain systems with state delay~\cite{7353149}.
Here, assume that the uncertainty $\Delta$ is known and
that $z_\Delta$, $w_\Delta=\Delta z_\Delta$, and $w_{\mrm{p}}$ are available
to the controller. For $\underline{\Psi}_\Delta\neq I$, find the realization of $\Psi_\Delta = \mbf{D}(\underline{\Psi}_\Delta)$,
\begin{equation*}
\bmat{
\mcl T_{{\Psi}_\Delta}
\dot{\mbf{x}}_{{\Psi}_\Delta}(t)\\
y(t)
}
=
\bmat{
\mcl A_{{\Psi}_\Delta}
&\mcl B_{{\Psi}_\Delta}^{\mrm{z}}
&\mcl B_{{\Psi}_\Delta}^{\mrm{w}}\\
I&0&0
}
\bmat{
\mbf{x}_{{\Psi}_\Delta}(t)\\
z_\Delta(t)\\
w_\Delta(t)
}.
\end{equation*}
We impose the zero initial condition
$\mbf{x}_{{\Psi}_\Delta}(0)=0$. The second block row gives
$y(t)=\mbf{x}_{{\Psi}_\Delta}(t)$. The complete controller has the PIE
realization
\begin{equation*}
\begin{aligned}
\bmat{
\mcl T_{{\Psi}_\Delta}
\dot{\mbf{x}}_{{\Psi}_\Delta}(t)\\
u(t)
}
&=
\bmat{
\mcl A_{{\Psi}_\Delta}
&0
&\mcl B_{{\Psi}_\Delta}^{\mrm{z}}
&\mcl B_{{\Psi}_\Delta}^{\mrm{w}}\\
\mcl K_\Psi&\mcl K_{\mrm{G}} & 0 & 0
}
\bmat{
\mbf{x}_{{\Psi}_\Delta}(t)\\
\mbf{x}(t)\\
z_\Delta(t)\\
w_\Delta(t)
}.
\end{aligned}
\end{equation*}
Hence $\mbf{x}_{{\Psi}_\Delta}$ is the controller state. The plant state
$\mbf{x}$ and the signals that enter the dynamic part of the filter are
external controller inputs, while $u$ is its output. This
realization makes explicit the signal-availability assumption stated in
Section~\ref{sec:primal-controller-analysis}.
